\documentclass[12pt,a4paper]{article}

\usepackage[margin=1in]{geometry}
\usepackage{amsmath, amssymb}
\usepackage{booktabs}
\usepackage{enumitem}
\usepackage{hyperref}
\usepackage{parskip}
\usepackage{titlesec}
\usepackage{array}
\usepackage{xcolor}
\usepackage{fancyhdr}
\setlength{\headheight}{14.5pt}
\addtolength{\topmargin}{-2.5pt}
\usepackage{graphicx}
\usepackage{rotating}
\usepackage{multirow}
\usepackage{longtable}
\usepackage{pdflscape}

\hypersetup{
    colorlinks=true,
    linkcolor=blue!60!black,
    urlcolor=blue!60!black,
}

\pagestyle{fancy}
\fancyhf{}
\rhead{Traceability Matrix}
\lhead{QUBO-Based DFL Node Placement}
\cfoot{\thepage}

\titleformat{\section}{\large\bfseries}{}{0em}{}[\titlerule]
\titleformat{\subsection}{\normalsize\bfseries}{}{0em}{}

\newcommand{\YES}{\textcolor{green!50!black}{\textbf{X}}}
\newcommand{\req}[1]{\textbf{#1}}

\title{
    \textbf{Traceability Matrix} \\[0.5em]
    \large QUBO-Based Sensor Node Placement for Device-Free Localization
}
\author{Mohamed Khalil Brik}
\date{Spring 2026}

\begin{document}

\maketitle
\tableofcontents
\newpage

% ─────────────────────────────────────────────────────────────────────────────
\section{Introduction}

This document provides the full traceability matrix for the QUBO-based
sensor node placement system. It links every functional requirement from
the System Requirements and Design Specifications to the implementation
component that fulfills it and the test case in the Verification and
Validation Report that verifies it.

The matrix gives reviewers a single reference point to confirm that no
requirement is left unimplemented or untested.

% ─────────────────────────────────────────────────────────────────────────────
\section{Document References}

\begin{table}[h!]
\centering
\renewcommand{\arraystretch}{1.3}
\begin{tabular}{@{}ll@{}}
\toprule
\textbf{Abbreviation} & \textbf{Document} \\
\midrule
SRS & System Requirements and Design Specifications \\
ESC & Engineering Standards and Constraints \\
VVR & Verification and Validation Report \\
\bottomrule
\end{tabular}
\end{table}

% ─────────────────────────────────────────────────────────────────────────────
\section{Requirement Descriptions}

The table below summarizes each functional requirement for quick reference.

\begin{table}[h!]
\centering
\renewcommand{\arraystretch}{1.3}
\begin{tabular}{@{}lp{10cm}@{}}
\toprule
\textbf{ID} & \textbf{Description} \\
\midrule
FR-1  & Load all input CSV files (RSS matrices, coordinates, pairs,
        importance, redundancy). \\
FR-2  & Build QUBO matrix $Q$ parameterized by $\alpha$ using importance
        and redundancy values. \\
FR-3  & Run binary search over $\alpha \in [0,1]$ with tolerance
        $\epsilon = 0.01$ to find a budget-feasible selection. \\
FR-4  & Solve the QUBO using SA (\texttt{SASampler}) and SQA
        (\texttt{SQASampler}) with specified read and sweep counts. \\
FR-5  & Implement a greedy baseline that selects pairs by descending
        importance while respecting the same AP and MP budgets. \\
FR-6  & Evaluate any selection using 3-NN localization and report the
        mean Euclidean error in meters over the test set. \\
FR-7  & Compute the Time-to-Solution (TTS) at $p_R = 0.99$ for SA and SQA. \\
FR-8  & Repeat the full binary search 100 times with varying seeds and
        report mean, standard deviation, min, and max statistics. \\
FR-9  & Save all numerical results to \texttt{data/output/} as CSV files. \\
FR-10 & Generate all ten publication-quality figures and save them to
        \texttt{figures/} at 300~DPI. \\
\bottomrule
\end{tabular}
\caption{Functional requirement summary.}
\end{table}

% ─────────────────────────────────────────────────────────────────────────────
\section{Requirements to Implementation Traceability}

\begin{table}[h!]
\centering
\renewcommand{\arraystretch}{1.4}
\begin{tabular}{@{}lp{8.5cm}l@{}}
\toprule
\textbf{Requirement} & \textbf{Implementing Component} & \textbf{File(s)} \\
\midrule
FR-1 &
    CSV loading with \texttt{pd.read\_csv}; shape printed to console
    for runtime verification. &
    \texttt{scripts/01}, \texttt{scripts/04} \\

FR-2 &
    \texttt{build\_qubo(alpha, I, R)} function constructs the sparse
    QUBO dictionary. &
    \texttt{scripts/01}, \texttt{scripts/04} \\

FR-3 &
    \texttt{binary\_search()} function; \texttt{while b - a >= EPSILON}
    loop with feasibility check. &
    \texttt{scripts/01}, \texttt{scripts/04} \\

FR-4 &
    \texttt{oj.SASampler().sample\_qubo()} with 100 reads / 1000 sweeps;
    \texttt{oj.SQASampler().sample\_qubo()} with 10 reads / 1000 sweeps. &
    \texttt{scripts/01}, \texttt{scripts/04} \\

FR-5 &
    Greedy loop over importance-sorted pairs; set-based AP/MP budget check. &
    \texttt{scripts/03}, \texttt{scripts/04} \\

FR-6 &
    \texttt{localizer()} / \texttt{knn\_error()} function; Euclidean
    distance in RSS space; mean error over test set. &
    \texttt{scripts/01}, \texttt{scripts/03}, \texttt{scripts/04} \\

FR-7 &
    \texttt{calculate\_tts(p\_s, tau, p\_target)} function;
    $\tau$ from wall-clock time (SA) or 20~$\mu$s (SQA). &
    \texttt{scripts/01}, \texttt{scripts/04} \\

FR-8 &
    Outer loop over \texttt{range(N\_RUNS)} with \texttt{seed = run};
    \texttt{pandas} groupby statistics saved to summary CSV. &
    \texttt{scripts/04} \\

FR-9 &
    \texttt{df.to\_csv()} calls at the end of each script save seven
    output files. &
    \texttt{scripts/01}, \texttt{scripts/03}, \texttt{scripts/04} \\

FR-10 &
    \texttt{matplotlib} figures with IEEE \texttt{rcParams};
    saved via \texttt{fig.savefig()}. &
    \texttt{scripts/02}, \texttt{scripts/05},
    \texttt{scripts/06}, \texttt{scripts/07} \\
\bottomrule
\end{tabular}
\caption{Requirements to implementation traceability.}
\end{table}

% ─────────────────────────────────────────────────────────────────────────────
\section{Requirements to Test Case Traceability}

\begin{table}[h!]
\centering
\renewcommand{\arraystretch}{1.4}
\begin{tabular}{@{}llp{7cm}@{}}
\toprule
\textbf{Requirement} & \textbf{Test Case(s)} & \textbf{What is Verified} \\
\midrule
FR-1  & TC-01 & Array shapes match expected dimensions. \\
FR-2  & TC-02 & QUBO diagonal and off-diagonal terms at $\alpha = 0$ and $\alpha = 1$. \\
FR-3  & TC-03, TC-04 & Search terminates within 7 iterations; all feasible solutions are within budget. \\
FR-4  & TC-03, TC-04 & Solver is called correctly; budget check applied to its output. \\
FR-5  & TC-08, TC-10 & Greedy is deterministic; QUBO accuracy is no worse than greedy. \\
FR-6  & TC-05, TC-06 & Correct error for full selection; infinity for empty selection. \\
FR-7  & TC-07 & Correct TTS at $p_s = 0$ and $p_s \geq 1$ boundary cases. \\
FR-8  & TC-11, TC-12, TC-14 & Statistics are consistent; row count is correct; budget holds in all 100 trials. \\
FR-9  & TC-09 & All seven output files exist and are non-empty. \\
FR-10 & TC-13 & All figures saved at 300~DPI. \\
\bottomrule
\end{tabular}
\caption{Requirements to test case traceability.}
\end{table}

% ─────────────────────────────────────────────────────────────────────────────
\section{Full Traceability Matrix}

The table below consolidates all three dimensions: requirements,
implementation components, and test cases.

\begin{landscape}
\begin{table}[p]
\centering
\renewcommand{\arraystretch}{1.4}
\small
\begin{tabular}{@{}l p{4.2cm} p{3.8cm} p{3.5cm} l@{}}
\toprule
\textbf{Req.} &
\textbf{Description (short)} &
\textbf{Implementing Script} &
\textbf{Key Function / Component} &
\textbf{Test Cases} \\
\midrule
FR-1  & Data ingestion
      & \texttt{01}, \texttt{04}
      & \texttt{pd.read\_csv}, shape checks
      & TC-01 \\

FR-2  & QUBO construction
      & \texttt{01}, \texttt{04}
      & \texttt{build\_qubo()}
      & TC-02 \\

FR-3  & Binary search
      & \texttt{01}, \texttt{04}
      & \texttt{binary\_search()} loop
      & TC-03, TC-04 \\

FR-4  & SA / SQA solvers
      & \texttt{01}, \texttt{04}
      & \texttt{SASampler}, \texttt{SQASampler}
      & TC-03, TC-04 \\

FR-5  & Greedy baseline
      & \texttt{03}, \texttt{04}
      & Importance-sorted greedy loop
      & TC-08, TC-10 \\

FR-6  & k-NN evaluation
      & \texttt{01}, \texttt{03}, \texttt{04}
      & \texttt{localizer()} / \texttt{knn\_error()}
      & TC-05, TC-06 \\

FR-7  & TTS computation
      & \texttt{01}, \texttt{04}
      & \texttt{calculate\_tts()}
      & TC-07 \\

FR-8  & Multi-run experiment
      & \texttt{04}
      & 100-run loop with seed = run
      & TC-11, TC-12, TC-14 \\

FR-9  & Result persistence
      & \texttt{01}, \texttt{03}, \texttt{04}
      & \texttt{df.to\_csv()} calls
      & TC-09 \\

FR-10 & Figure generation
      & \texttt{02}, \texttt{05}, \texttt{06}, \texttt{07}
      & \texttt{matplotlib} with IEEE rcParams
      & TC-13 \\
\bottomrule
\end{tabular}
\caption{Full requirements--implementation--test traceability matrix.}
\end{table}
\end{landscape}

% ─────────────────────────────────────────────────────────────────────────────
\section{Coverage Summary}

\begin{itemize}[leftmargin=2em]
    \item \textbf{Requirements coverage:} 10 of 10 functional requirements
          are implemented (100\%).
    \item \textbf{Test coverage:} 10 of 10 functional requirements are
          covered by at least one test case (100\%).
    \item \textbf{Test pass rate:} 14 of 14 test cases pass (100\%).
    \item \textbf{NFR coverage:} All five non-functional requirements
          (NFR-1 through NFR-5) are addressed in the Engineering Standards
          and Constraints document and verified through code inspection.
\end{itemize}

\end{document}
